%!Mode:: "TeX:UTF-8"
%!TEX TS-program = xelatex %!Mode:: "TeX:UTF-8"

\documentclass[a4paper]{rukh-cv}

\begin{document}
%
%----------------------------------------------------------------------------------------
%	LEFT COLUMN
%----------------------------------------------------------------------------------------
%
\begin{minipage}[t]{0.20\textwidth} % The left column takes up 33% of the text width of the page

\namesection{Peng}{Ding}

\infosection{(+86)18611105508 \\
\href{mailto:i@dingpeng.me}{i@dingpeng.me}  \\ 
Room 1-309 \\ 
FIT Building \\ 
Tsinghua Campus \\ 
Haidian, Beijing}{\includegraphics[width=.9em]{icons/linkedin.pdf} \href{https://www.linkedin.com/in/dingpeng}{dingpeng} \\ 
\includegraphics[width=.9em]{icons/github.pdf} \href{https://github.com/rxkh}{rxkh}}

\end{minipage} % The end of the left column
\hfill
%
%----------------------------------------------------------------------------------------
%	RIGHT COLUMN
%----------------------------------------------------------------------------------------
%
\begin{minipage}[t]{0.74\textwidth} % The right column takes up 66% of the text width of the page

%------------------------------------------------
% Education
%------------------------------------------------

\section{Education}

\runsubsection{Tsinghua University}
\descript{| MEng in Software Engineering}
\location{Expected Sep 2013 - Jul 2016 | Beijing, China}
\vspace{\topsep}
\begin{tightitemize}
\item Study on topic modeling and trend prediction.
\end{tightitemize}

\sectionspace

%------------------------------------------------

\runsubsection{Tsinghua University}
\descript{| BEng in Computer Software}
\location{Sep 2009 - Jul 2013 | Beijing, China}
\begin{tightitemize}
\item Outstanding student leaders of Tsinghua University.
\item Merit scholarship on social work of Tsinghua University.
\end{tightitemize}

\sectionspace

%------------------------------------------------
% Projects
%------------------------------------------------

\section{Projects}

\runsubsection{CodingYoung}
\descript{| Public Coding Course}
\location{Apr 2014 – Jun 2015 | Beijing, China}
\begin{tightitemize}
\item Developed a 12-hour coding course \emph{Python Programming in Action}.
\item The course has been selected as an elective course in a key middle school, and has helped 25 students age 13 to 15 learn to code in Python.
\end{tightitemize}

\sectionspace

%------------------------------------------------

\runsubsection{Marketing Solutions for ANSYS China}
\descript{| Enterprise Apps}
\location{Apr 2014 – Jan 2015 | Beijing, China}
\begin{tightitemize}
\item Implemented the new version of official WeChat subscription account, whose amount of subscribers increased 6 times in the first 6 months.
\item Implemented a monitor tool to crawl contents on the official website (en-us) and provide weekly updates reports.
\end{tightitemize}

\sectionspace

%------------------------------------------------

\runsubsection{Student-Tsinghua}
\descript{| Campus Website}
\location{Aug 2013 – Jul 2014 | Beijing, China}
\begin{tightitemize}
\item A website for extracurricular activity management in Tsinghua University, which has 20000 individual users and 1700 PV in average.
\item In charge of product design, technical architecture and schedule control of new features such as activity resources management, new front page, etc.
\end{tightitemize}

\sectionspace

%------------------------------------------------

\runsubsection{MICAPS}
\descript{| Weather Forecast Production System}
\location{Jan 2013 – Feb 2013 | Beijing, China}
\begin{tightitemize}
\item The system is applied in National Meteorological Center of China, and it is used for authoritative weather forecast production.
\item Designed an interactive meteorological-editing algorithm based on isoline dragging and re-implemented the related module,
    which decreased the running time on each interaction from about 10s down to 0.1s.
\end{tightitemize}

\sectionspace

%------------------------------------------------

\runsubsection{Course Rating Website}
\descript{| Social Network Web App}
\location{May 2011 – Jan 2013 | Beijing, China}
\begin{tightitemize}
\item It is a web app built on Renren Open Platform that allows students at Tsinghua University to rate, comment and share experience about school courses.
\item Got 5,000 users in first 2 weeks and 15,000 in total, ranked 2nd among educational apps on Renren.com for 2 weeks.
\end{tightitemize}

\sectionspace

%------------------------------------------------
% Experience
%------------------------------------------------

\section{Experience}

\runsubsection{Beijing Baijugu Technology Co., Ltd.}
\descript{| Co-founder}
\location{Apr 2014 – Apr 2015 | Beijing, China}
\begin{tightitemize}
\item Implementation of epibolic websites, mobile app servers, WeChat account services, etc.
\item Design and implementation of education automation tools for high schools.
\end{tightitemize}

\sectionspace

%------------------------------------------------

\runsubsection{Youth League Committee, Tsinghua University}
\descript{| Director}
\location{Aug 2013 – Jul 2014 | Beijing, China}
\begin{tightitemize}
\item Led a tech team to build and maintain the Studeng-Tsinghua website.
\end{tightitemize}

\sectionspace

%------------------------------------------------
% Skills
%------------------------------------------------

\section{Skills}

Python, Django, C/C++, C\#, Java, SQL, MongoDB

\sectionspace

\end{minipage}

\end{document}
